\documentclass[11pt,a4paper]{article}
\usepackage[T1]{fontenc}
\usepackage[utf8]{inputenc}
\usepackage[main=italian,provide=*]{babel}
\usepackage{amsmath,amssymb}
\usepackage{geometry}
\geometry{margin=2.3cm,top=2.5cm,bottom=2.7cm}
\usepackage{booktabs}
\usepackage[table,dvipsnames]{xcolor}
\usepackage{tcolorbox}
\tcbuselibrary{skins,breakable}
\usepackage{titlesec}
\usepackage{enumitem}
\usepackage{needspace}
\usepackage{fancyhdr}
\usepackage{graphicx}
\usepackage{hyperref}
\hypersetup{colorlinks=true,linkcolor=Sepia,citecolor=Sepia,urlcolor=Sepia}
\newcommand{\Tr}{\operatorname{Tr}}
\newcommand{\ket}[1]{\lvert #1 \rangle}
\newcommand{\bra}[1]{\langle #1 \rvert}

% --- palette ---
\definecolor{inkblue}{HTML}{1B3A5C}
\definecolor{lightblue}{HTML}{EAF1F8}
\definecolor{accent}{HTML}{B0562A}
\definecolor{softgray}{HTML}{6B6B6B}
\definecolor{okgreen}{HTML}{2E6B3E}
\definecolor{okgreenbg}{HTML}{EAF5EC}
\definecolor{warnred}{HTML}{9E3B3B}
\definecolor{warnredbg}{HTML}{FBEDED}

% --- titoli di sezione ---
\titleformat{\section}
  {\normalfont\Large\bfseries\color{inkblue}}
  {\thesection}{0.6em}{}[{\color{inkblue!40}\titlerule}]
\titlespacing{\section}{0pt}{0.9em}{0.4em}
\titleformat{\subsection}
  {\normalfont\large\bfseries\color{accent}}
  {}{0em}{}
\titlespacing{\subsection}{0pt}{0.6em}{0.2em}

\pagestyle{fancy}
\fancyhf{}
\renewcommand{\headrulewidth}{0.4pt}
\fancyhead[L]{\small\color{softgray}VQE noise-aware sul dimero --- resoconto completo}
\fancyhead[R]{\small\color{softgray}\thepage}
\fancyfoot[C]{\small\color{softgray}Tesi triennale, Samuele --- Relatore: Prof.\ Alessandro Chiesa}

% --- box riutilizzabili ---
\newtcolorbox{keyeq}[1][]{
  colback=lightblue, colframe=inkblue!70, boxrule=0.6pt,
  arc=2pt, left=8pt, right=8pt, top=3pt, bottom=3pt, #1
}
\newtcolorbox{okbox}[1]{
  colback=okgreenbg, colframe=okgreen!60, boxrule=0.6pt,
  arc=2pt, left=7pt, right=7pt, top=2pt, bottom=2pt,
  fonttitle=\bfseries\color{okgreen}, title={#1}
}
\newtcolorbox{warnbox}[1]{
  colback=warnredbg, colframe=warnred!60, boxrule=0.6pt,
  arc=2pt, left=7pt, right=7pt, top=2pt, bottom=2pt,
  fonttitle=\bfseries\color{warnred}, title={#1}
}

\title{\vspace{-1.5em}\color{inkblue}VQE noise-aware sul dimero\\[0.2em]
\Large\normalfont\color{softgray}Resoconto completo: processo, risultati, considerazioni, conclusioni}
\author{Note di lavoro --- Tesi triennale, Samuele\\ Relatore: Prof.\ Alessandro Chiesa}
\date{}

\begin{document}
\sloppy
\maketitle
\thispagestyle{fancy}
\vspace{-2.0em}

File di codice: \texttt{vqe\_noise\_aware\_dimero.py},
\texttt{validate\_vqe\_noise\_aware\_dimero.py}. Punto di lavoro:
$b/J=0.35$, $D/J=0.80$, $J=1$ (``test 2''). Modello di rumore:
\texttt{noise\_model\_dimero.py} (modello di rumore, calibrazione reale
\texttt{ibm\_torino}, arXiv:2504.15187).

\tableofcontents

\section{Obiettivo e collocazione}

Il VQE con termine DM sotto rumore, stadio della pipeline di rumore
(Parte 2, già chiusa in tutti e
cinque i passi) misura la \emph{degradazione} di un risultato ideale:
l'ansatz PMA-2q$\cdot$3 è ottimizzato \textbf{senza} rumore in Parte 1,
poi il circuito con quella ricetta viene eseguito su un simulatore
\textbf{con} rumore acceso. Questo documento affronta la domanda
complementare, lasciata esplicitamente aperta in \texttt{log\_decisioni.md}:

\begin{keyeq}
\centering
Se il rumore è acceso già durante l'ottimizzazione variazionale --- come
avverrebbe su un hardware reale, dove non esiste una versione ideale a
cui accedere per calibrare l'ottimizzatore --- l'ottimizzatore converge a
parametri diversi? Il risultato finale (energia, fedeltà) cambia?
\end{keyeq}

Questa è un'estensione \textbf{facoltativa} del VQE con termine DM
sotto rumore, non uno step
mancante della pipeline richiesta dal relatore. Lo scope è dichiarato
esplicitamente e limitato a tre parti:
\begin{enumerate}[leftmargin=1.4em,itemsep=0.2em]
\item rioptimizzare \textbf{solo} il VQE con termine DM sotto rumore;
\item un \textbf{solo} controllo a valle: verificare se questo sposta
$N^*=8$, il risultato centrale della quantum simulation (Trotter)
sotto rumore;
\item i correlatori dinamici sotto rumore ($N^*=5$) e lo scan sui
parametri di rumore \textbf{non} vengono ritoccati qui --- l'estensione
a questi due stadi si appoggia a un argomento di struttura, non a una riverifica
numerica diretta (criterio di proporzionalità, vista la scadenza di
laurea di settembre 2026).
\end{enumerate}

\section{Argomento fisico e previsione}

Il canale di rumore usato (costruito nel modello di rumore) è
depolarizzante, isotropo: riduce
uniformemente la sfera di Bloch in ogni direzione, senza una direzione
privilegiata. Per un canale di questo tipo, l'azione su uno stato puro
$\ket{\psi(\theta)}$ preparato dall'ansatz si scrive schematicamente

\begin{keyeq}
\centering
$\rho(\theta) = (1-\lambda)\,\ket{\psi(\theta)}\bra{\psi(\theta)} + \lambda\,\dfrac{\mathbf{1}}{d}$
\end{keyeq}

con $\lambda>0$ legato a $\varepsilon_{1q},\varepsilon_{2q}$ tramite le
relazioni derivate in \texttt{teoria\_modello\_rumore.tex}, e $d=4$ per 2
qubit. L'energia rumorosa diventa allora

\begin{keyeq}
\centering
$E_\text{rumoroso}(\theta) = \Tr[\rho(\theta) H] = (1-\lambda)\,E_\text{ideale}(\theta) + \lambda\,\dfrac{\Tr[H]}{d}$
\end{keyeq}

cioè una trasformazione \textbf{affine} di $E_\text{ideale}(\theta)$ in
funzione di $\theta$: il secondo termine non dipende da $\theta$, il
primo è $E_\text{ideale}(\theta)$ moltiplicato per una costante positiva
$1-\lambda$. Una trasformazione affine con coefficiente positivo non
cambia la posizione del minimo: \textbf{lo stesso $\theta^*$ che minimizza
$E_\text{ideale}$ minimizza anche $E_\text{rumoroso}$}. A questo si
aggiunge una seconda osservazione, verificata indipendentemente (§4.1):
il conteggio CNOT del circuito transpilato non dipende dai valori
numerici dei parametri, quindi non c'è nemmeno una via indiretta (via
costo in gate) per cui convenga spostarsi verso un $\theta$ diverso da
quello ideale.

\begin{keyeq}[title={Previsione dichiarata prima della verifica}]
\centering
$\theta^*_\text{rumore} \approx \theta^*_\text{ideale}$, nessun vantaggio
misurabile a rioptimizzare sotto rumore rispetto al semplice riuso dei
parametri ideali.
\end{keyeq}

Questa è una previsione, non un fatto acquisito: il resto del documento
la mette alla prova numericamente.

\subsection*{Perché un solo controllo a valle basta per i correlatori dinamici e lo scan sui parametri}

La quantum simulation (Trotter) sotto rumore ha già dimostrato che il
contributo della preparazione VQE è
una costante additiva indipendente da $N$ nella fedeltà $F(N)$: non entra
né nel termine che decresce con $N$ (errore di Trotter, $\sim a/N$) né in
quello che cresce con $N$ (rumore accumulato, $\sim g\varepsilon_{2q}N$).
I correlatori dinamici e lo scan sui parametri condividono la stessa
struttura additiva e la stessa preparazione VQE, quindi lo stesso
argomento si applica identico a tutti e tre gli stadi a valle: se il
controllo su $N^*$ (quantum simulation Trotter sotto rumore) conferma che
la preparazione noise-aware non cambia nulla, la stessa conclusione vale
per i correlatori dinamici e lo scan sui parametri senza bisogno di
ripetere il controllo lì. Se invece lo
smentisse, si riaprirebbero anche quelli, perché la premessa che li
chiudeva sarebbe caduta --- per questo il controllo numerico resta
comunque necessario, non solo l'argomento.

\section{Implementazione}

\subsection{Formulazione del problema di ottimizzazione}

L'ansatz è lo stesso PMA-2q$\cdot$3 di Parte 1 (blocco RBS-conservante il
numero di eccitazioni, più due rotazioni $R_y$ indipendenti, 3 parametri).
L'energia rumorosa, per un dato $\theta$, si calcola simulando il
circuito su \texttt{AerSimulator(method="density\_matrix")} con il
\texttt{NoiseModel} del modello di rumore agganciato:

\begin{keyeq}
\centering
$E_\text{rumoroso}(\theta) = \Tr[\rho(\theta) H], \qquad \rho(\theta) = \mathcal{N}\!\big(\ket{\psi(\theta)}\bra{\psi(\theta)}\big)$
\end{keyeq}

dove $\mathcal N$ è il canale completo (depolarizzante 1q/2q + errore di
lettura, se applicabile). Poiché la simulazione a operatore densità è
\textbf{esatta} (nessun campionamento a shot finiti), $E_\text{rumoroso}(\theta)$
è liscia in $\theta$: gli stessi ottimizzatori del caso ideale (COBYLA
multistart + polish L-BFGS-B, \texttt{vqe\_test2.py}) restano applicabili
senza modifiche concettuali.

\subsection{Una trappola di transpilazione scoperta e corretta}

Il primo tentativo di implementazione applicava all'ottimizzazione lo
stesso principio già validato per il blocco di Trotter ripetuto nella
quantum simulation Trotter sotto rumore: transpilare l'ansatz
\textbf{una sola volta}, in forma parametrica
(parametri simbolici, non ancora assegnati), e poi riusare
\texttt{assign\_parameters} ad ogni valutazione dell'obiettivo --- per
evitare di ritranspilare a ogni chiamata durante l'ottimizzazione.

\begin{warnbox}{Verificato: qui questo principio non si applica}
Con parametri simbolici, il transpilatore non può fondere le sequenze di
rotazioni a 1 qubit il cui angolo dipende dal \emph{valore numerico} di
$\theta$ (a differenza del blocco di Trotter, dove il blocco è
letteralmente identico ad ogni ripetizione e non contiene parametri
variazionali da ottimizzare). Il risultato è un circuito strutturalmente
più rumoroso: transpilando \emph{prima} di assegnare i parametri si
ottengono 19 \texttt{rz} + 14 \texttt{sx} contro gli 11 \texttt{rz} + 8
\texttt{sx} che si ottengono transpilando \emph{dopo}, a parità di 2
CNOT (Fig.~\ref{fig:bug}).
\end{warnbox}

\begin{figure}[h]
\centering
\includegraphics[width=0.62\linewidth]{figure/fig_bug_conteggio_gate.pdf}
\caption{\textit{Conteggio dei gate nel circuito transpilato, a parità di
parametri $\theta$, a seconda dell'ordine fra transpilazione e
assegnazione dei parametri. Il conteggio CNOT resta invariato (2), ma i
gate a singolo qubit quasi raddoppiano se si transpila prima
dell'assegnazione: ogni gate rz/sx in più è un'ulteriore occasione di
errore per il canale depolarizzante.}}
\label{fig:bug}
\end{figure}

Conseguenza pratica osservata durante lo sviluppo: il VQE noise-aware con
questa implementazione convergeva a un'energia $E\approx-3.4753$,
\textbf{peggiore} del semplice riuso dei parametri ideali sotto rumore
($E=-3.50164605$, versione originale del VQE con termine DM sotto
rumore). Questo è di per sé un segnale
inequivocabile di errore implementativo, non un risultato fisico: un
ottimizzatore che valuta l'obiettivo rumoroso vero non può in linea di
principio fare peggio di un punto ammissibile già noto ($\theta^*_\text{ideale}$
è comunque un punto del dominio di ottimizzazione). La correzione
--- transpilare \textbf{dopo} l'assegnazione dei parametri, ad ogni
valutazione, coerentemente con \texttt{vqe\_dm\_rumoroso\_dimero.py} ---
elimina la discrepanza (§4).

Costo computazionale della scelta corretta: $\sim9\,\text{ms}$ per
valutazione (transpilazione + simulazione a operatore densità),
trascurabile rispetto al numero di valutazioni tipico di un multistart
COBYLA ($R\times\text{maxiter}$ dell'ordine delle migliaia).

\section{Verifica}

\subsection{Invarianza del conteggio CNOT rispetto a \texorpdfstring{$\theta$}{theta}}

Premessa fisica della previsione (§2): il conteggio CNOT del circuito
transpilato (nell'ordine corretto, dopo l'assegnazione) non deve
dipendere dai valori numerici di $\theta$. Verificato su 5 assegnazioni
casuali indipendenti:

\begin{center}
\renewcommand{\arraystretch}{1.2}
\begin{tabular}{@{}ccc@{}}
\toprule
\textbf{Prova} & \textbf{CNOT} & \textbf{Dettaglio gate}\\
\midrule
1--5 & $2$ (tutte) & $11$ \texttt{rz} $+\ 8$ \texttt{sx} $+\ 2$ \texttt{cx} (tutte)\\
\bottomrule
\end{tabular}
\end{center}

\noindent CNOT $=2$ in tutte le prove, conteggio dei gate a 1 qubit anch'esso
costante in questo campione: nessuna direzione del paesaggio dei
parametri verso cui convenga spostarsi per ``risparmiare'' rumore.

\subsection{Limite di rumore nullo}

Il VQE noise-aware, eseguito con il \texttt{NoiseModel} spento, deve
ritrovare l'ottimo esatto già noto da Parte 1:

\begin{center}
\renewcommand{\arraystretch}{1.2}
\begin{tabular}{@{}lc@{}}
\toprule
\textbf{Quantità} & \textbf{Valore}\\
\midrule
$E_0$ (esatto, registrato Parte 1) & $-3.57321145$\\
$E$ (noise-aware, rumore nullo) & $-3.57321145$\\
$|\Delta E|$ & $1.32\times10^{-12}$\\
\bottomrule
\end{tabular}
\end{center}

\subsection{Risultato principale: confronto energia e fedeltà sotto rumore}

Risultato con il rumore di riferimento (\texttt{ibm\_torino}), multistart
$R=6$ punti casuali $+\ 1$ seme uguale a $\theta^*_\text{ideale}$ (per
garantire di non fare mai peggio del semplice riuso, vista la dispersione
discussa in §4.4):

\begin{center}
\renewcommand{\arraystretch}{1.25}
\begin{tabular}{@{}lcc@{}}
\toprule
\textbf{Quantità} & \textbf{Riuso $\theta^*_\text{ideale}$ (VQE+DM rumoroso)} & \textbf{VQE noise-aware}\\
\midrule
$E$ (rumoroso) & $-3.50164605$ & $-3.50164616$\\
$F$ (rumorosa) & $0.98500601$ & $0.98500602$\\
CNOT & $2$ & $2$\\
\bottomrule
\end{tabular}
\end{center}

\begin{center}
\renewcommand{\arraystretch}{1.2}
\begin{tabular}{@{}cc@{}}
\toprule
\textbf{Quantità} & \textbf{Valore}\\
\midrule
$\Delta E$ (noise-aware $-$ riuso ideale) & $-1.16\times10^{-7}$\\
$\Delta F$ (noise-aware $-$ riuso ideale) & $+1.76\times10^{-8}$\\
\bottomrule
\end{tabular}
\end{center}

\noindent Scostamento sei ordini di grandezza più piccolo della scala
fisica del problema (perdita di fedeltà $\sim1.5\times10^{-2}$ a rumore
acceso): conferma diretta dell'argomento (§2), non solo entro la
tolleranza attesa ma sostanzialmente a precisione numerica.

Lo stesso confronto, ripetuto sull'intero scan di $\varepsilon_{2q}$ già
usato nella verifica della versione originale, mostra le due curve
sovrapposte su tutto il range esplorato (Fig.~\ref{fig:eps2q}):

\begin{figure}[h]
\centering
\includegraphics[width=0.95\linewidth]{figure/fig_scan_eps2q_confronto.pdf}
\caption{\textit{Energia e fedeltà rumorose in funzione di
$\varepsilon_{2q}$ (a $\varepsilon_{1q}=2.9\times10^{-4}$,
$p_\text{readout}=2.3\times10^{-2}$ fissati), confronto fra il riuso dei
parametri ideali e la rioptimizzazione noise-aware. Le due curve sono
indistinguibili a questa scala su tutto il range esplorato, dal rumore
nullo a $\varepsilon_{2q}=1.5\times10^{-2}$ ($4\times$ il valore di
riferimento).}}
\label{fig:eps2q}
\end{figure}

\subsection{Robustezza del multistart: un'osservazione onesta}

Un controllo diverso da quello atteso inizialmente: con $R=6$ punti di
partenza casuali e \textbf{nessun} seme (stesso numero di restart
della versione originale, applicato qui al VQE noise-aware), i risultati
dipendono dal seme dell'RNG. Su 8 semi diversi, lo scarto dall'ottimo
noto ($E=-3.50164616$) varia da $3.6\times10^{-6}$ (il migliore) a
$6.7\times10^{-3}$ (il peggiore) --- una dispersione reale, dovuta al
numero finito di restart in un paesaggio energetico con più minimi
locali vicini fra loro, \textbf{non} imputabile al bug di transpilazione
descritto in §3.2 (già corretto in questi test).

\begin{figure}[h]
\centering
\includegraphics[width=0.8\linewidth]{figure/fig_robustezza_multistart.pdf}
\caption{\textit{Energie finali dei singoli restart COBYLA (rumore
acceso), per tre configurazioni di multistart: $R=6$ senza seme (6
punti), $R=6$+seme (7 punti, il seme è $\theta^*_\text{ideale}$),
$R=20$ senza seme (20 punti). La linea tratteggiata è l'ottimo noto. Con
soli 6 restart casuali il minimo trovato può restare a
$\sim10^{-3}$--$10^{-2}$ dall'ottimo; con 20 restart il punto migliore si
avvicina all'ottimo entro $3.6\times10^{-6}$.}}
\label{fig:multistart}
\end{figure}

\begin{okbox}{Il seme non altera la conclusione fisica}
Includere $\theta^*_\text{ideale}$ come uno dei punti di partenza del
multistart garantisce che il risultato non sia mai peggiore del semplice
riuso (per costruzione: è un punto ammissibile già valutato). Un
multistart indipendente, più ampio ($R=20$, nessun seme), converge allo
\textbf{stesso} ottimo rumoroso trovato con $R=6$+seme entro
$2.6\times10^{-6}$ --- conferma che il risultato principale non è un
artefatto del seme, ma una proprietà genuina del paesaggio energetico,
coerente con l'argomento (§2) secondo cui $\theta^*_\text{rumore}\approx\theta^*_\text{ideale}$.
\end{okbox}

\subsection{Controllo a valle: \texorpdfstring{$N^*$}{N*} della quantum simulation (Trotter) sotto rumore}

Ultimo controllo richiesto dallo scope (§1, punto 2): rieseguire lo scan
di fedeltà della quantum simulation Trotter sotto rumore
(\texttt{trotter\_rumoroso\_dimero.py}) usando
$\theta^*_\text{noise-aware}$ al posto di $\theta^*_\text{ideale}$ come
preparazione, per vedere se il massimo si sposta.

\begin{center}
\renewcommand{\arraystretch}{1.2}
\begin{tabular}{@{}lcc@{}}
\toprule
\textbf{Parametri usati come preparazione} & $\boldsymbol{N^*}$ & $\boldsymbol{F(N^*)}$\\
\midrule
$\theta^*_\text{ideale}$ (versione originale, Trotter rumoroso) & $8$ & $0.817725$\\
$\theta^*_\text{noise-aware}$ (qui) & $8$ & $0.817767$\\
\bottomrule
\end{tabular}
\end{center}

\begin{figure}[h]
\centering
\includegraphics[width=0.95\linewidth]{figure/fig_F_vs_N_confronto.pdf}
\caption{\textit{Sinistra: fedeltà $F(N)$ rispetto allo stato bersaglio
fisico, rumore di riferimento acceso, con le due preparazioni a
confronto --- il massimo resta a $N^*=8$ in entrambi i casi. Destra:
scostamento $F_\text{noise-aware}-F_\text{ideale}$ in funzione di $N$,
scala ingrandita ($\sim10^{-4}$): la differenza è visibile solo
ingrandendo di quattro ordini di grandezza rispetto alla scala del
pannello di sinistra, ed è più pronunciata a $N$ piccolo (dove la
preparazione pesa relativamente di più sul numero totale di gate) che a
$N$ grande.}}
\label{fig:FvsN}
\end{figure}

\begin{okbox}{$N^*$ non si sposta}
Il massimo di $F(N)$ resta a $N^*=8$, con lo stesso valore di fedeltà
entro $4\times10^{-5}$. La previsione (b)--(c) di §2 è quindi confermata
numericamente, non solo argomentata: il contributo di preparazione resta
il termine costante (floor) già caratterizzato nella quantum
simulation Trotter sotto rumore, indipendente
da \emph{quale} versione dei parametri VQE (ideali o noise-aware) lo
genera.
\end{okbox}

\section{Considerazioni}

\begin{itemize}[leftmargin=1.3em,itemsep=0.35em]
\item \textbf{La previsione teorica (§2) è confermata, non solo
plausibile.} L'argomento della trasformazione affine non era ovvio a
priori che si traducesse in uno scostamento numerico trascurabile (poteva
esistere un effetto del second'ordine, per esempio dovuto al fatto che il
canale a 2 qubit non è esattamente descritto dal modello semplificato
usato nella derivazione schematica di §2). Il controllo numerico mostra
che l'effetto, se esiste, è ben al di sotto della precisione rilevante
per la tesi ($\Delta E\sim10^{-7}$, sei ordini di grandezza sotto la
scala fisica del problema).

\item \textbf{Il bug scoperto (§3.2) è un caso utile da riportare in
tesi.} È l'immagine speculare della ``trappola di transpilazione'' già
nota per la quantum simulation Trotter sotto rumore (ritranspilare
l'intero circuito di Trotter fonde gli
$N$ blocchi, cancellando la dipendenza da $N$): lì il problema era
transpilare troppo tardi (dopo aver composto tutti i blocchi); qui è
transpilare troppo presto (prima di conoscere i valori numerici dei
parametri variazionali). Il principio comune: \emph{la scelta di quando
transpilare rispetto a quando si fissano i valori numerici non è neutra,
va sempre verificata esplicitamente per ogni nuovo contesto d'uso, non
assunta per analogia da un contesto precedente}.

\item \textbf{Sul ruolo del seme nel multistart.} Includere un punto di
partenza noto (qui: l'ottimo ideale di Parte 1) è una scelta pragmatica
legittima quando quel punto è comunque disponibile, ma non deve
mascherare che il paesaggio energetico rumoroso ha minimi locali
raggiungibili con un numero insufficiente di restart casuali (Fig.~
\ref{fig:multistart}). In un'esecuzione reale su hardware, senza accesso
a una versione ideale precalcolata, servirebbe un $R$ più grande (o
un'euristica di inizializzazione diversa) per raggiungere con
confidenza lo stesso ottimo.

\item \textbf{Limite dell'argomento (§2).} La derivazione della
trasformazione affine è schematica (assume che il canale depolarizzante
a 2 qubit produca esattamente una miscela lineare stato-ideale/stato
massimamente misto); il modello reale include anche l'errore di lettura
e il fatto che il rumore agisce gate-per-gate lungo il circuito, non come
un singolo canale applicato all'intero stato finale. La conferma
numerica (§4.3) è quindi più forte della sola derivazione euristica, e va
citata come tale in tesi: il risultato è verificato, non solo previsto.
\end{itemize}

\section{Conclusioni ed estensione ai correlatori dinamici e allo scan sui parametri}

Il controllo numerico richiesto dallo scope conferma la previsione fisica
di §2: per un canale depolarizzante isotropo, rioptimizzare l'ansatz
PMA-2q$\cdot$3 sotto rumore non cambia in modo misurabile né l'energia/
fedeltà del VQE+DM rumoroso ($\Delta E\sim10^{-7}$) né la posizione di
$N^*$ della quantum simulation Trotter sotto rumore (invariato, $N^*=8$).
Per l'argomento di struttura (§2.1), la stessa conclusione si estende ai
correlatori dinamici sotto rumore ($N^*=5$) e allo scan sui parametri di
rumore, che condividono l'identica preparazione VQE
e l'identica struttura additiva
$\varepsilon_\text{tot}(N)\simeq a/N+g\varepsilon_{2q}N+\delta_\text{prep}$:
non riverificati numericamente qui per il criterio di proporzionalità
dichiarato in §1, ma la premessa che li giustificava --- il floor di
preparazione è costante e indipendente da $N$ --- non solo non è stata
falsificata da questo controllo, è stata rinforzata.

\begin{okbox}{Sintesi in una riga}
Estensione facoltativa CHIUSA: rioptimizzare sotto rumore non serve a
nulla di misurabile in questo problema, e questo è stato dimostrato, non
solo assunto.
\end{okbox}

\needspace{6\baselineskip}
\section*{File prodotti}
\addcontentsline{toc}{section}{File prodotti}

\texttt{vqe\_noise\_aware\_dimero.py} (implementazione, funzioni di
ottimizzazione e di controllo), \texttt{validate\_vqe\_noise\_aware\_dimero.py}
(i cinque controlli di §4), questo documento. Aggiornamento corrispondente
in \texttt{log\_decisioni.md} (sessione 5 settembre 2026).

\end{document}
